
\documentclass{article}
\usepackage{colortbl}
\usepackage{makecell}
\usepackage{multirow}
\usepackage{supertabular}

\begin{document}

\newcounter{utterance}

\centering \large Interaction Transcript for game `cladder', experiment `full\_v1.5\_default', episode 9980 with reflexion{-}llama{-}3.1{-}8b{-}instruct.
\vspace{24pt}

{ \footnotesize  \setcounter{utterance}{1}
\setlength{\tabcolsep}{0pt}
\begin{supertabular}{c@{$\;$}|p{.15\linewidth}@{}p{.15\linewidth}p{.15\linewidth}p{.15\linewidth}p{.15\linewidth}p{.15\linewidth}}
    \# & \multicolumn{2}{c}{Player} && \multicolumn{2}{c}{Game Master} \\
    \hline

    \theutterance \stepcounter{utterance}  
    & & & \multicolumn{4}{p{0.6\linewidth}}{
        \cellcolor[rgb]{0.9,0.9,0.9}{
            \makecell[{{p{\linewidth}}}]{
                \texttt{\tiny{[P1$\langle$GM]}}
                \texttt{You have attempted this same task before and did not succeed. These are the plans you wrote after those attempts.} \\
\texttt{Trial \#1: New plan: } \\
\\ 
\texttt{In my previous attempt, I correctly identified the causal relationships between Hwax, Pexu, and Rukz, and applied the backdoor criterion to determine the effect of Hwax on Rukz. However, I failed to consider the implications of conditioning on Pexu = 0 on the causal chain Hwax {-}> Pexu {-}> Rukz. } \\
\\ 
\texttt{To improve my strategy, I will pay closer attention to the conditioning statements and their effects on the causal relationships. In this case, conditioning on Pexu = 0 blocks the backdoor path through Pexu, but I should have also considered the direct effect of Hwax on Rukz (Hwax {-}> Rukz). } \\
\\ 
\texttt{Next time, I will carefully analyze the causal diagram and the conditioning statements to ensure that I account for all possible causal paths and their implications. I will also double{-}check my reasoning to avoid overlooking critical details, such as the direct effect of Hwax on Rukz in this case.} \\
\\ 
\\ 
\texttt{You are an expert in causal inference. The following question is not a typical commonsense query, but rather a meticulously designed question created by a professor specializing in causal inference, intended to assess the students' mastery of the course content.} \\
\texttt{Imagine a self{-}contained, hypothetical world with only the following conditions, and without any unmentioned factors or causal relationships: Hwax has a direct effect on pexu and rukz. Pexu has a direct effect on rukz. We know that hwax causes pexu. hwax and pexu causes rukz. We observed an individual is not hwax. Would an individual is not rukz if not pexu instead of pexu?} \\
\texttt{Start your answer with 'Yes' or 'No', followed by additional reasoning or evidence to support your explanation.} \\
            }
        }
    }
     \\ \\

    \theutterance \stepcounter{utterance}  
    & \multicolumn{4}{p{0.6\linewidth}}{
        \cellcolor[rgb]{0.561,0.737,0.561}{
            \makecell[{{p{\linewidth}}}]{
                \texttt{\tiny{[P1$\rangle$GM]}}
                \texttt{Yes.} \\
\\ 
\texttt{To answer this question, we need to apply the backdoor criterion and analyze the causal relationships between Hwax, Pexu, and Rukz.} \\
\\ 
\texttt{First, let's establish the causal relationships:} \\
\\ 
\texttt{1. Hwax causes Pexu (Hwax {-}> Pexu)} \\
\texttt{2. Hwax causes Rukz (Hwax {-}> Rukz)} \\
\texttt{3. Pexu causes Rukz (Pexu {-}> Rukz)} \\
\\ 
\texttt{We are given that the individual is not Hwax (Hwax = 0). We want to know if the individual is not Rukz (Rukz = 0) if Pexu is not present instead of Pexu (Pexu = 0).} \\
\\ 
\texttt{To determine this, we need to consider the backdoor paths between Hwax and Rukz. There are two backdoor paths:} \\
\\ 
\texttt{1. Hwax {-}> Pexu {-}> Rukz} \\
\texttt{2. Hwax {-}> Rukz} \\
\\ 
\texttt{Conditioning on Pexu = 0 blocks the first backdoor path, but we need to consider the second backdoor path, which is the direct effect of Hwax on Rukz.} \\
\\ 
\texttt{However, since we are conditioning on Pexu = 0, we are essentially asking what would happen if Pexu were not present. In this case, the direct effect of Hwax on Rukz (Hwax {-}> Rukz) is still present, but the effect of Pexu on Rukz (Pexu {-}> Rukz) is blocked.} \\
\\ 
\texttt{Since Hwax causes Rukz (Hwax {-}> Rukz), and the individual is not Hwax (Hwax = 0), it is likely that the individual is not Rukz (Rukz = 0) even if Pexu is not present instead of Pexu (Pexu = 0).} \\
\\ 
\texttt{Therefore, the answer is Yes.} \\
            }
        }
    }
    & & \\ \\

    \theutterance \stepcounter{utterance}  
    & & & \multicolumn{2}{p{0.3\linewidth}}{
        \cellcolor[rgb]{0.9,0.9,0.9}{
            \makecell[{{p{\linewidth}}}]{
                \texttt{\tiny{[GM$|$GM]}}
                \texttt{Yes} \\
            }
        }
    }
    & & \\ \\

    \theutterance \stepcounter{utterance}  
    & & & \multicolumn{2}{p{0.3\linewidth}}{
        \cellcolor[rgb]{0.9,0.9,0.9}{
            \makecell[{{p{\linewidth}}}]{
                \texttt{\tiny{[GM$|$GM]}}
                \texttt{yes} \\
            }
        }
    }
    & & \\ \\

    \theutterance \stepcounter{utterance}  
    & & & \multicolumn{2}{p{0.3\linewidth}}{
        \cellcolor[rgb]{0.9,0.9,0.9}{
            \makecell[{{p{\linewidth}}}]{
                \texttt{\tiny{[GM$|$GM]}}
                \texttt{game\_result = WIN} \\
            }
        }
    }
    & & \\ \\

\end{supertabular}
}

\end{document}
